% SHARD-ID: AISM-63-TOPOLOGY-TOOLBOX
% SHARD-TITLE: The algebraic-topology toolbox: quotient manifolds, triangulation, Lefschetz--Hopf, local index sign, orientable top cohomology, Kunneth, and Hopf structure
% SHARD-SUMMARY: Reproduces lem-topology-quotient-manifold, the quotient-manifold theorem for smooth free proper Lie-group actions.
% SHARD-SUMMARY: Reproduces lem-topology-finite-triangulation, the finite-triangulation consequence for compact smooth manifolds without boundary. Reproduces lem-topology-lefschetz-hopf, the maximal-simplex form of the Lefschetz--Hopf fixed-point formula. Reproduces lem-topology-local-index-sign, the determinant-sign formula for a nondegenerate local fixed-point index.
% SHARD-SUMMARY: Reproduces lem-topology-orientable-top-cohomology, the nonvanishing of real top cohomology for a closed connected orientable manifold. Reproduces lem-topology-kunneth-cross-product, the cohomological Kunneth ring isomorphism and its finite-CW real-field specialization. Reproduces lem-topology-hopf-structure, the finite-dimensional exterior-algebra corollary of the characteristic-zero Hopf structure theorem.
% SHARD-KEYWORDS: algebraic topology, quotient manifold, finite triangulation, Lefschetz--Hopf, local fixed-point index, orientability, top cohomology, Kunneth theorem, Hopf algebra, af validated
\section{The algebraic-topology toolbox}\label{sec:topology-toolbox}

This section collects the seven \texttt{af}-validated topology rows used by the
Stage-1 fixed-point and cohomological arguments.  Each account reports the live proof
tree, including the source and scope repairs recorded during adversarial verification.

\begin{lemma}[\texttt{lem-topology-quotient-manifold}]\label{lem:topology-quotient-manifold}
Let a Lie group $G$ act smoothly, freely and properly on a smooth manifold $M$.
Then $M/G$ is a topological manifold of dimension $\dim M-\dim G$, and it has a
unique smooth structure for which the quotient map $\pi:M\to M/G$ is a smooth
submersion.
\end{lemma}

\noindent\textbf{Registry contract (verbatim).}
\contractquote{Quotient manifold theorem: if a Lie group G acts smoothly, freely, and properly on a smooth manifold M, then the orbit space M/G is a topological manifold of dimension dim M - dim G with a unique smooth structure for which the quotient map pi:M->M/G is a smooth submersion.}

\paragraph{Proof (prose account of the af-validated tree).}
The tree is a direct application of Lee's Quotient Manifold Theorem.
Its registered external has precisely the smoothness, freeness and properness
hypotheses above and returns the orbit-space dimension, the unique smooth structure
and the submersion property in one step.  The only elaboration is documentary:
the extracted source writes the orbit notation, minus sign and quotient-map symbol
with OCR control glyphs.  The validated children identify those glyphs with $M/G$,
$\dim M-\dim G$ and $\pi:M\to M/G$, while leaving the byte-matched external
unchanged.  Three verifier challenges observed that merely naming the theorem did
not yet place it in the formal claim context.  The repaired leaf therefore invokes
the registered citation alias explicitly and records why the interface's empty
optional context display does not cancel that binding.  With the external formally
in scope, its conclusion is exactly the registry contract.
The resulting four-node tree is taint clean, and the root adds no new assertion.
This is precisely the direct-source leaf recorded in the export.
No further geometric argument or additional premise is used.

\begin{lemma}[\texttt{lem-topology-finite-triangulation}]\label{lem:topology-finite-triangulation}
Every compact smooth manifold without boundary is homeomorphic to the realization of
a finite simplicial complex.
\end{lemma}

\noindent\textbf{Registry contract (verbatim).}
\contractquote{Finite triangulation of compact smooth manifolds: every compact smooth manifold without boundary is homeomorphic to a finite simplicial complex.}

\paragraph{Proof (prose account of the af-validated tree).}
Let $M$ satisfy the hypotheses.
Smoothness implies the $C^1$ regularity used in the
registered Munkres theorem, while the pinned terminology external identifies absence
of boundary with ``non-bounded.''  Munkres therefore supplies a simplicial complex
$K$ and a $C^1$ triangulation $f:|K|\to M$.  The two registered definition excerpts
make the important next step literal: such a triangulation is a homeomorphism onto
$M$.  Hence $|K|$ is compact, by applying the continuous inverse of $f$ to compact
$M$.  For each vertex $v$, the live tree takes the union of the relative interiors of
the simplices containing $v$; these open stars cover $|K|$, and a vertex $w$ belongs
to the star centered at $v$ exactly when $w=v$.  A finite subcover therefore captures
every vertex among finitely many centers.  There are then only finitely many simplices,
since each is a finite subset of the finite vertex set.  Thus $K$ is finite.
The validated workspace is the clean re-seeded Munkres route: the earlier Cairns route
was retired because its missing Veblen--Whitehead definitions could not be grounded
locally, and a later corrupted-locus re-seed was discarded rather than inherited.

\begin{lemma}[\texttt{lem-topology-lefschetz-hopf}]\label{lem:topology-lefschetz-hopf}
Let $f:X\to X$ be a map of a finite polyhedron with finitely many fixed points,
each lying in a maximal simplex of $X$.  Then
\[
  L(f)=\sum_{x\in\operatorname{Fix}(f)}\operatorname{ind}(f,x).
\]
\end{lemma}

\noindent\textbf{Registry contract (verbatim).}
\contractquote{Lefschetz-Hopf formula (maximal-simplex form): if f:X->X is a map of a finite polyhedron with a finite set of fixed points, each of which lies in a maximal simplex of X, then the Lefschetz number L(f) is the sum of the indices of all the fixed points of f.}

\paragraph{Proof (prose account of the af-validated tree).}
The two-node tree applies Arkowitz--Brown, Theorem~1.2, verbatim.
The registered external assumes a self-map of a finite polyhedron, a finite fixed
set, and placement of every fixed point in a maximal simplex.  Its conclusion is
exactly the displayed sum of the local indices, with the notation of
Definition~\ref{def:lefschetz-fixed-point-data}.
The proof tree changes no hypothesis and introduces no intermediate construction.
In particular, the maximal-simplex condition is retained: the earlier design
wording with only isolated fixed points was broader than the pinned theorem and
was not validated.
The sole child quotes the registered source statement literally.
The root then takes that child by direct theorem application.
There is no deformation, approximation or index comparison inside this row.
The validated conclusion is therefore the exact pinned source statement.
Thus any downstream use must establish maximal-simplex placement separately;
this row does not infer it from finiteness or isolation.

\begin{lemma}[\texttt{lem-topology-local-index-sign}]\label{lem:topology-local-index-sign}
Let $x$ be an isolated fixed point of a smooth self-map $f$ of a compact orientable
manifold.  If $\det(I-Df_x)\ne0$, then
\[
  \operatorname{ind}(f,x)=\sgn\det(I-Df_x).
\]
\end{lemma}

\noindent\textbf{Registry contract (verbatim).}
\contractquote{Nondegenerate local fixed-point index: if x is an isolated fixed point of a smooth self-map f of a compact orientable manifold with det(I-Df_x) != 0, then its local fixed-point index is sgn det(I-Df_x).}

\paragraph{Proof (prose account of the af-validated tree).}
The decisive live leaf applies the nondegenerate clause of
Definition~\ref{def:lefschetz-fixed-point-data} directly: the smooth compact orientable
scope and the determinant hypothesis are exactly its registered hypotheses, and $1$
and $I$ denote the same tangent-space identity.  The remaining live branches audit the
classical local calculation.  Writing $A=Df_x$, invertibility of $I-A$ excludes the
eigenvalue $1$.  After shrinking one chart ball so that its closed ball remains in the
coordinate domain, the Leray--Schauder formula reduces $J(f,x)$ to $J(A,0)$.
For $h=I-A$, the Brouwer-degree theorem gives
$d(h,V)=\sgn\det(I-A)$.  Independently, the linear Leray--Schauder formula gives
$J(A,0)=(-1)^\beta$, where $\beta$ counts characteristic values in $(0,1)$.
Reciprocation identifies these with real eigenvalues of $A$ larger than $1$;
conjugate nonreal pairs contribute positively to $\det(I-A)$, so
$(-1)^\beta=\sgn\det(I-A)$.  Verification repaired nine challenges concerning chart
compactness, an out-of-scope use on a noncompact vector space, cutoff inequalities and
unsupported locality bridges.  The final scope was user-ratified from the original
unqualified $C^1$ wording to the smooth compact orientable contract quoted above.

\begin{lemma}[\texttt{lem-topology-orientable-top-cohomology}]\label{lem:topology-orientable-top-cohomology}
If $M$ is a connected compact orientable $d$-manifold without boundary, then
$H^d(M;\mathbb R)\ne0$.
\end{lemma}

\noindent\textbf{Registry contract (verbatim).}
\contractquote{Top cohomology of a closed orientable manifold: if M is a connected compact orientable d-manifold without boundary, then H^d(M;R) != 0.}

\paragraph{Proof (prose account of the af-validated tree).}
Compactness without boundary is the ``closed'' hypothesis in the registered Hatcher
theorem.  The live tree first obtains real orientability without importing an
unregistered coefficient-change result.  Hatcher's registered top-cohomology
corollary supplies a nonzero class in $H^d(M;\mathbb R)$; the field-coefficient
universal coefficient theorem then forces $H_d(M;\mathbb R)\ne0$.  If $M$ were not
$\mathbb R$-orientable, Theorem~3.26(b) would inject this group into $\mathbb R$ with
image $\{r:2r=0\}=\{0\}$, a contradiction.  Theorem~3.26(a) now gives
$H_d(M;\mathbb R)\cong\mathbb R$.  Applying the same universal coefficient theorem
yields
\[
  H^d(M;\mathbb R)\cong
  \operatorname{Hom}_{\mathbb R}(H_d(M;\mathbb R),\mathbb R),
\]
and the right-hand side contains any chosen isomorphism
$H_d(M;\mathbb R)\to\mathbb R$, hence is nonzero.  A verifier challenge rejected the
tree's first change-of-coefficients shortcut; the preceding Corollary--UCT--Theorem
route is the validated replacement.

\begin{lemma}[\texttt{lem-topology-kunneth-cross-product}]\label{lem:topology-kunneth-cross-product}
Let $R$ be a commutative coefficient ring and let $X,Y$ be CW complexes such that
each $H^k(Y;R)$ is a finitely generated free $R$-module.  Then the cross product is a
ring isomorphism
\[
  H^*(X;R)\otimes_R H^*(Y;R)\xrightarrow{\ \cong\ }H^*(X\times Y;R).
\]
In particular, this holds over $R=\mathbb R$ for finite-CW spaces with
finite-dimensional cohomology.
\end{lemma}

\noindent\textbf{Registry contract (verbatim).}
\contractquote{Cohomological Kunneth isomorphism over R: for CW complexes X and Y with each H^k(Y;R) a finitely generated free R-module, the cross product H*(X;R) (x)_R H*(Y;R) -> H*(X x Y;R) is an isomorphism of rings; in particular this holds over the field R = reals for finite-CW spaces with finite-dimensional cohomology.}

\paragraph{Proof (prose account of the af-validated tree).}
Hatcher's registered Theorem~3.15 states the first assertion for CW complexes over
a commutative coefficient ring when every $H^k(Y;R)$ is finitely generated and free.
The live tree applies that theorem directly.  The commutative-ring wording records
the scope correction forced by verification, since the provisional ``arbitrary
ring'' reading was broader than the source.
For a noncommutative ring the displayed tensor product of two left modules need
not even be available in the asserted form.
For the special case, a finite-CW space is in particular a CW complex.
Over the field $\mathbb R$, every finite-dimensional
cohomology group has a finite basis, hence is a finitely generated free
$\mathbb R$-module.  These are exactly the hypotheses of the general clause,
so its cross product gives the stated ring isomorphism for the finite-CW pair.
The final tree node is this specialization by modus ponens.
No finiteness conclusion beyond those explicit hypotheses, and no assertion for
a noncommutative coefficient ring, is made.

\begin{lemma}[\texttt{lem-topology-hopf-structure}]\label{lem:topology-hopf-structure}
Every finite-dimensional connected graded-commutative bialgebra over a field of
characteristic zero is an exterior algebra on odd-degree homogeneous generators,
where ``connected graded'' uses the nonnegative grading convention.
\end{lemma}

\noindent\textbf{Registry contract (verbatim).}
\contractquote{Hopf structure theorem in the form consumed by Stage 1: a finite-dimensional connected graded-commutative bialgebra over a characteristic-zero field is an exterior algebra on odd-degree homogeneous generators.}

\paragraph{Proof (prose account of the af-validated tree).}
Let $A=\bigoplus_{n\ge0}A_n$ be the bialgebra over the characteristic-zero field $F$.
Connectedness gives $A_0=F1$.  For homogeneous $a\in A_n$, $n>0$, degree preservation
and the counit identities force the coproduct decomposition
\[
  \Delta a=a\otimes1+1\otimes a+\sum_i a_i'\otimes a_i'',
  \qquad |a_i'|,|a_i''|>0.
\]
Thus the hypotheses match Hatcher's registered Hopf-algebra convention, and total
finite-dimensionality makes every graded piece finite-dimensional.  The structure
theorem gives an algebra isomorphism
$A\cong\Lambda_F(V_{\mathrm{odd}})\otimes_F F[V_{\mathrm{even}}]$.
If $V_{\mathrm{even}}$ contained a polynomial generator $x$, the independent monomials
$1,x,x^2,\ldots$ would make its polynomial algebra infinite-dimensional.  The map
$p\mapsto1\otimes p$ is injective, split by the exterior augmentation, so the tensor
product and hence $A$ would also be infinite-dimensional.  Therefore the polynomial
factor has no generators and $A\cong\Lambda_F(V_{\mathrm{odd}})$.  Verification
repaired one scope gap by making explicit that ``connected graded'' carries the
source's nonnegative convention; negatively graded counterexamples are not claimed.
